\documentclass[conference]{IEEEtran}
\usepackage{cite}
\usepackage{amsmath,amssymb,amsfonts}
\usepackage{algorithmic}
\usepackage{graphicx}
\usepackage{textcomp}
\usepackage{xcolor}
\usepackage{booktabs}
\usepackage{float}
\usepackage{hyperref}
\usepackage{multirow}

\def\BibTeX{{\rm B\kern-.05em{\sc i\kern-.025em b}\kern-.08em
    T\kern-.1667em\lower.7ex\hbox{E}\kern-.125emX}}

\title{3D Volumetric Detection of Pulmonary Nodules in CT Scans using Advanced 3D U-Net}

\author{\IEEEauthorblockN{Do Hong Chung}
\IEEEauthorblockA{\textit{Department of Data Science} \\
\textit{University of Science and Technology of Hanoi}\\
Hanoi, Vietnam \\
chungdh.22ba13055@usth.edu.vn}
\IEEEauthorblockA{Student ID: 22BA13055}
}

\begin{document}

\maketitle

\begin{abstract}
The early detection of pulmonary nodules in 3D Computed Tomography (CT) scans is a critical step in diagnosing lung cancer. In this study, I address the task of 3D volumetric semantic segmentation to identify spherical nodule lesions. Utilizing the LUNA16 dataset, I propose a patch-based deep learning pipeline driven by an Advanced 3D U-Net architecture. To overcome the extreme class imbalance inherent in 3D medical volumes and computational constraints, the model operates on localized $64 \times 64 \times 64$ voxel patches and optimizes a combined Dice and Binary Cross Entropy (BCE) loss function. Experimental results over 30 training epochs demonstrate that the model successfully converges, achieving a 3D Dice score of 0.4885 and localizing 3D nodule structures effectively.
\end{abstract}

\begin{IEEEkeywords}
Pulmonary Nodule, 3D Semantic Segmentation, 3D U-Net, LUNA16, Medical Imaging, Volumetric Detection.
\end{IEEEkeywords}

\section{Introduction}
Lung cancer remains one of the leading causes of cancer-related mortality worldwide. Computed Tomography (CT) scans provide detailed 3D representations of the lungs, making them the gold standard for detecting early-stage pulmonary nodules. However, manually inspecting hundreds of cross-sectional slices per patient is exceptionally time-consuming for radiologists.

The objective of this practice is to build a 3D Machine Learning model to automatically segment these volumetric nodules. Unlike 2D imaging (such as X-rays), 3D CT data presents immense computational challenges. A single scan can consume gigabytes of memory, and the target nodules often occupy less than 1\% of the entire lung volume. I focus on implementing a patch-based extraction strategy and an Advanced 3D U-Net to tackle these spatial and hardware constraints.

\section{Dataset and Preprocessing}

\subsection{Dataset Description}
I utilized a subset (\texttt{subset0}) of the \textbf{LUNA16 (LUng Nodule Analysis 2016)} dataset, which is a benchmark for nodule detection.
\begin{itemize}
    \item \textbf{Input:} 3D CT Scans stored in \texttt{.mhd} and \texttt{.raw} formats.
    \item \textbf{Annotations:} A \texttt{CSV} file providing the real-world coordinates (X, Y, Z in millimeters) and the diameter of each confirmed nodule.
\end{itemize}

\subsection{3D Preprocessing Pipeline}
To ensure the model fits within GPU memory constraints (such as an NVIDIA T4) and trains efficiently, several crucial preprocessing steps were applied:
\begin{itemize}
    \item \textbf{Coordinate Transformation:} Nodule coordinates were converted from physical World coordinates (mm) to continuous Voxel indices using the origin and spacing metadata of each scan.
    \item \textbf{Patch Extraction:} Instead of loading the entire 3D volume, $64 \times 64 \times 64$ spatial patches were cropped around the center of each nodule.
    \item \textbf{Hounsfield Unit (HU) Windowing:} The radiodensity values were clipped to $[-1000, 400]$ to isolate lung tissues and bone structures, followed by min-max normalization to $[0, 1]$.
    \item \textbf{Mask Generation:} Ground truth 3D spherical masks were dynamically generated based on the nodule's physical diameter converted into voxel space.
\end{itemize}

\section{Methodology}

\subsection{Model Architecture: Advanced 3D U-Net}
To capture the complex volumetric context of pulmonary vessels and nodules, an \textbf{Advanced 3D U-Net} was designed.
\begin{itemize}
    \item \textbf{Encoder (Feature Extraction):} Utilizes 3D Convolutional layers ($3 \times 3 \times 3$ kernels). To increase the receptive field, the network depth was extended to four levels, with feature channels doubling at each step ($16 \rightarrow 32 \rightarrow 64 \rightarrow 128$).
    \item \textbf{Decoder (Spatial Recovery):} Employs trilinear upsampling combined with skip connections from the encoder to reconstruct the 3D spatial resolution accurately.
    \item \textbf{Output:} A final $1 \times 1 \times 1$ Convolution with a Sigmoid activation outputs a 3D probability map indicating the likelihood of a voxel belonging to a nodule.
\end{itemize}

\subsection{Loss Function and Training Configuration}
In a $64^3$ volume, nodule voxels are severely outnumbered by healthy background voxels. Using standard BCE Loss would lead to a trivial solution where the model predicts only the background. To counter this, a \textbf{Combined Loss Function} was implemented:
$$ \mathcal{L}_{total} = \mathcal{L}_{BCE} + \mathcal{L}_{Dice} $$
The model was trained using the Adam optimizer ($lr=2e-4$) for 30 epochs and a batch size of 4. Furthermore, 3D spatial augmentations (random flipping across X, Y, and Z axes) were applied to prevent overfitting.

\section{Results}

\subsection{Qualitative Analysis}
Because visualizing a 3D volume on a 2D paper is difficult, Fig. \ref{fig:res} presents the middle axial cross-section (Slice Z=32) of the $64 \times 64 \times 64$ volume. The AI prediction successfully localized the nodule's center, effectively ignoring surrounding noise and healthy tissues.

\begin{figure}[htbp]
\centering
\includegraphics[width=\linewidth]{Visualization.png}
\caption{Visualization of 3D Volumetric Segmentation (Cross-section at Z=32). Left to Right: Input CT Slice, Ground Truth Spherical Mask, and Predicted Nodule Region.}
\label{fig:res}
\end{figure}

\subsection{Quantitative Evaluation}
The model's performance was evaluated using the 3D Volumetric Dice Coefficient and 3D Intersection over Union (IoU). The results are summarized in Table \ref{tab1}.

\begin{table}[htbp]
\caption{Performance Metrics on 3D Patches (30 Epochs)}
\begin{center}
\begin{tabular}{|c|c|c|}
\hline
\textbf{Target Class} & \textbf{Mean 3D Dice Score} & \textbf{Mean 3D IoU} \\
\hline
\textbf{Pulmonary Nodule} & \textbf{0.4885} & \textbf{0.3821} \\
\hline
\end{tabular}
\label{tab1}
\end{center}
\end{table}

\subsection{Training Convergence}
Fig. \ref{fig:loss} illustrates the optimization progress over 30 epochs. 

\begin{figure}[htbp]
\centering
\includegraphics[width=\linewidth]{Training-loss-curve.png}
\caption{Training Loss Curve combining BCE and Dice Loss (30 Epochs).}
\label{fig:loss}
\end{figure}

The combined BCE+Dice loss started at approximately \textbf{1.59} in the initial epoch. Driven by the Adam optimizer, the loss exhibited a remarkably smooth and consistent descent, eventually converging to slightly below \textbf{1.20} by the 30th epoch. This steady curve indicates that the Advanced 3D U-Net effectively learned the complex spatial features of nodules without suffering from vanishing gradients or severe overfitting.

\section{Discussion}
Transitioning from 2D to 3D segmentation introduces significant hurdles, which is reflected in the metrics.
\begin{itemize}
    \item \textbf{Metric Sensitivity:} A 3D Dice score of \textbf{0.4885} and an IoU of \textbf{0.3821} are highly realistic for this task. Achieving a high 3D Dice score is geometrically much more difficult than in 2D. A minor offset of just 1-2 voxels at the boundary of a tiny 3D sphere drastically reduces the intersection volume. Thus, these obtained scores prove the model's solid capability for preliminary small-lesion detection.
    \item \textbf{False Positives:} In 3D space, blood vessels can often mimic the spherical appearance of nodules when viewed in isolated patches. The extended depth of the Advanced U-Net (up to 128 filters) helped provide the necessary contextual awareness to distinguish between continuous vessels and isolated nodules, ensuring a steady training convergence.
\end{itemize}

\section{Conclusion}
This study successfully established a robust 3D deep learning pipeline for detecting pulmonary nodules from raw CT scans. By utilizing a patch-based approach, combining BCE and Dice Loss, and leveraging an Advanced 3D U-Net, the computational bottleneck was resolved while maintaining meaningful localization accuracy (Dice 0.4885). Future improvements could involve sliding-window inference to scan entire, uncropped lung volumes and deploying 3D Attention mechanisms to further refine the boundaries of detected nodules.

\begin{thebibliography}{00}
\bibitem{b1} O. Ronneberger, P. Fischer, and T. Brox, "U-Net: Convolutional Networks for Biomedical Image Segmentation," in \textit{MICCAI}, 2015.
\bibitem{b2} Özgün Çiçek et al., "3D U-Net: Learning Dense Volumetric Segmentation from Sparse Annotation," in \textit{MICCAI}, 2016.
\bibitem{b3} LUNA16 Grand Challenge, "LUng Nodule Analysis 2016," \url{https://luna16.grand-challenge.org/}.
\end{thebibliography}

\end{document}